\documentclass{article}
\usepackage[utf8]{inputenc}
\usepackage[polish]{babel}
\usepackage{booktabs}
\usepackage{tabularx}
\usepackage{geometry}
\geometry{a4paper, margin=1in}

\begin{document}

\section*{Plan Studiów: Matematyka Teoretyczna (I stopień)}

\noindent \textbf{Status na dzień:} 26 lutego 2026 r. (Semestr letni 2025/2026)

\begin{table}[h!]
\centering
\small
\begin{tabularx}{\textwidth}{@{} l X c l @{}}
\toprule
\textbf{Okres} & \textbf{Przedmioty i kamienie milowe} & \textbf{ECTS} & \textbf{Status} \\ \midrule
\textbf{Semestry 1-2} & Analiza 1-3, Algebra Liniowa 1-2, Wstęp do mat., Algebra 1R, & 88 & W realizacji \\ 
(2025/2026) & Analiza numeryczna, Programowanie 1, Topologia, Teoria zbiorów & & \\ \midrule
\textbf{Przerwa letnia} & \textbf{Samoedukacja:} Równania różniczkowe 1R, Rachunek prawd. 1R & (14) & Przygotowanie \\ \midrule
\textbf{Semestr 3} & Analiza funkcjonalna 1, Miara i całka, Proseminarium 1, & 32 & Zaplanowane \\ 
(Zimowy 26/27) & Projekt, Numerical methods (T), Przedmiot T, Ochrona WI, WF & & \\ \midrule
\textbf{Semestr 4} & RR 1R, RP 1R, Rozmaitości różniczkowalne, Funkcje anal. R, & 49* & Zaplanowane \\ 
(Letni 26/27) & Proseminarium 2, Algebra abstrakcyjna 2R (T), Szeregi (T), AF 2 (T) & & \\ \midrule
\textbf{Rok 3} & Przedmioty humanistyczne, Lektorat B2 (dopełnienie), & $\sim$11+ & Do ustalenia \\ 
(2027/2028) & \textbf{Egzamin Dyplomowy} & & \\ \bottomrule
\end{tabularx}
\caption*{* Faktyczne obciążenie w semestrze 4. spadnie do \textbf{35 ECTS} dzięki wakacyjnej nauce RR i RP.}
\end{table}

\section*{Weryfikacja wymogów specjalności teoretycznej}

\begin{itemize}
    \item \textbf{Przedmioty obowiązkowe (Tabela 7):} Zrealizujesz 100\% wymaganych kursów (Topologia, Miara i całka, Funkcje anal., Rozmaitości, AF 1, Proseminaria).
    \item \textbf{Pula przedmiotów „T” (min. 24 ECTS):} Twój plan przewiduje \textbf{32 ECTS} (Numerical Methods, Algebra 2R, Szeregi, AF 2, Wybrany przedmiot T), co spełnia wymóg z zapasem.
    \item \textbf{Punkty ECTS:} Po 4. semestrze będziesz mieć ok. 169 pkt. Do uzyskania dyplomu wymagane jest \textbf{180 ECTS}[cite: 29, 164].
\end{itemize}

\end{document}
